\documentclass[11pt,a4paper]{article}
\usepackage[T1]{fontenc}
\usepackage[utf8]{inputenc}
\usepackage[czech]{babel}
\usepackage{lmodern,microtype}
\usepackage{amsmath,amssymb,mathtools}
\usepackage{geometry}
\geometry{left=19mm,right=19mm,top=21mm,bottom=21mm}
\usepackage[dvipsnames]{xcolor}
\usepackage{enumitem,tabularx,booktabs}
\usepackage[most]{tcolorbox}
\usepackage{fancyhdr,lastpage}
\usepackage[hidelinks]{hyperref}
\definecolor{deep}{HTML}{203955}
\definecolor{teal}{HTML}{176E76}
\definecolor{slate}{HTML}{506173}
\definecolor{pale}{HTML}{F3F7F9}
\setlength{\parindent}{0pt}
\setlength{\parskip}{5pt}
\setlist[enumerate]{leftmargin=6mm,itemsep=5pt,topsep=3pt}
\pagestyle{fancy}\fancyhf{}
\fancyhead[L]{\sffamily\footnotesize\color{slate} MAT · 4. ročník · kombinatorika}
\fancyhead[R]{\sffamily\footnotesize\color{slate} Hodina \LESSONNUM\, / 112}
\fancyfoot[L]{\sffamily\footnotesize\color{slate} Výukový podklad · samostatný dokument}
\fancyfoot[R]{\sffamily\footnotesize\color{slate} \thepage/\pageref{LastPage}}
\renewcommand{\headrulewidth}{0.35pt}
\renewcommand{\footrulewidth}{0pt}
\newcommand{\sect}[1]{\par\vspace{2mm}{\Large\sffamily\bfseries\color{deep} #1}\par\vspace{-0.5mm}\noindent\color{teal}\rule{\linewidth}{0.7pt}\color{black}\vspace{1mm}}
\tcbset{theobox/.style={enhanced,breakable,colback=pale,colframe=teal!45!white,boxrule=.6pt,arc=1mm,left=3mm,right=3mm,top=2mm,bottom=2mm}}
\newcommand{\LESSONNUM}{36}
\begin{document}

{\Huge\sffamily\bfseries\color{deep} Kombinační číslo: výpočet a rovnice}\par
{\sffamily\large\color{teal} Výukový podklad · 36. vyučovací hodina}\par
\textbf{Hodinová dotace:} 45 minut\quad \textbf{Ročník:} 4. SPŠ\quad \textbf{Rozpis:} TP 2.2\par
\begin{tcolorbox}[theobox]\textbf{Výukový cíl.} Počítat kombinační čísla bez přetečení faktoriálů, používat krajní hodnoty a řešit přístupné rovnice.\end{tcolorbox}
\sect{Teorie a vysvětlení}
Kombinační číslo \(\binom nk\) čteme „\(n\) nad \(k\)“ a pro \(n\in\mathbb N_0\), \(0\le k\le n\) definujeme
\[\binom nk=\frac{n!}{k!(n-k)!}.\]
Zejména \(\binom n0=\binom nn=1\), \(\binom n1=n\), \(\binom n2=\frac{n(n-1)}2\).\par
\textbf{Výpočetní strategie:} zkrať faktoriály dříve, než budeš násobit. Pro \(k\) blízké \(n\) je vhodné vybrat menší z hodnot \(k\) a \(n-k\); proč to funguje, dokážeme v další hodině.\par
\textbf{Rovnice:} při \(\binom n2=45\) se řeší kvadratická rovnice s podmínkou, že \(n\) je celé a \(n\ge2\).
\sect{Řešené příklady}
\textbf{Příklad 1.} Vypočti \(\binom{10}{3}\).\par
\textit{Řešení.} \(\frac{10!}{3!7!}=\frac{10\cdot9\cdot8}{3\cdot2\cdot1}=120\).\par\smallskip
\textbf{Příklad 2.} Vypočti \(\binom{12}{10}\) bez velkých faktoriálů.\par
\textit{Řešení.} \(\binom{12}{10}=\frac{12!}{10!2!}=\frac{12\cdot11}{2}=66\).\par\smallskip
\textbf{Příklad 3.} Urči \(n\), jestliže \(\binom n2=45\).\par
\textit{Řešení.} \(\frac{n(n-1)}2=45\), tedy \(n^2-n-90=0=(n-10)(n+9)\). Přípustné je pouze \(n=10\).\par\smallskip
\textbf{Příklad 4.} Vyřeš \(\binom n2=3n\) pro \(n\ge2\).\par
\textit{Řešení.} \(n(n-1)/2=3n\). Protože \(n>0\), vydělíme \(n\): \(n-1=6\). Dostaneme \(n=7\).\par\smallskip
\textbf{Příklad 5.} Z osmi lidí se vybírá čtyřčlenná komise. Dva konkrétní kandidáti musejí oba být vybráni.\par
\textit{Řešení.} Dva jsou již předepsáni, z ostatních šesti vybereme ještě dva: \(\binom62=15\).\par\smallskip
\clearpage\sect{Procvičování pro studenty}
\begin{enumerate}
\item Spočti \(\binom83\).
\item Spočti \(\binom{11}{9}\).
\item Najdi \(n\), pro něž \(\binom n1=13\).
\item Najdi \(n\ge2\), pro které \(\binom n2=28\).
\item Ze sedmi prvků vyber sedm, případně žádný.
\end{enumerate}
\sect{Řešení a zdůvodnění k procvičování}
\begin{enumerate}
\item \(8\cdot7\cdot6/(3\cdot2\cdot1)=56\).
\item \(11\cdot10/2=55\).
\item Protože \(\binom n1=n\), \(n=13\).
\item \(n(n-1)=56\), tedy \(n=8\).
\item V obou případech existuje právě jeden výběr: \(\binom77=\binom70=1\).
\end{enumerate}
\sect{Metodické poznámky}
\textbf{Typické chyby:}
\begin{itemize}[leftmargin=5mm,itemsep=1pt]
\item Kombinační číslo je celé číslo; výpočet ve faktoriálech zkrať před vyčíslením.
\item Nikdy nezapomenout \(0!\) a krajní hodnoty \(k=0,n\).
\item Kořeny rovnic musí splnit podmínky \(n\in\mathbb N_0\) a \(k\le n\).
\end{itemize}
\textbf{Doporučený průběh 45 minut:} 0–8 definice; 8–20 zkracování; 20–32 rovnice; 32–42 cvičení; 42–45 závěr.\par
\textbf{Návaznost:} Vlastnosti kombinačních čísel vysvětlíme nejen algebraicky, ale i počítáním skupin dvěma způsoby.\par
\vspace{2mm}\noindent{\footnotesize\color{slate} \textbf{Východiska:} rozpis výuky MAT, 4. ročník, hodiny 29--41; tematický plán SPŠ; srovnání tematického rozsahu a náročnosti s realisticky.cz, oddíl Kombinatorika (lekce 090101--090119). Výklad, příklady a řešení jsou nově autorsky sestavené, nikoli převzatá sbírka.}\par
\end{document}